% File: lezione_01.tex
\chapter{Lezione 1}
\label{chap:lezione_01} % Etichetta per riferirsi al capitolo

\begin{flushright}
\textit{Data: 01/10/2025}
\end{flushright}


\section{Invarianza di Scala}

Uno dei concetti fondamentali nella fisica dei sistemi complessi è l'invarianza di scala. Questa proprietà descrive sistemi le cui caratteristiche statistiche non cambiano sotto una trasformazione di scala, ovvero un ingrandimento o una riduzione (zoom).

Consideriamo una trasformazione di scala sulla variabile $\epsilon$:
\begin{equation}
\epsilon \rightarrow \epsilon' = r\epsilon
\end{equation}
dove $r$ è il fattore di scala. Un sistema presenta invarianza di scala se una sua proprietà $S(\epsilon)$ si trasforma come segue:
\begin{equation}
S(\epsilon') = S(r\epsilon) = A(r)S(\epsilon)
\end{equation}
dove $A(r)$ è una funzione che dipende solo dal fattore di scala $r$ e non dalla scala originale $\epsilon$.

Si può dimostrare matematicamente che l'unica funzione che soddisfa questa relazione è la legge di potenza (power-law):
\begin{equation}
S(\epsilon) = a\epsilon^k
\end{equation}
dove $a$ e $k$ sono costanti.

Un esempio classico è il problema della misurazione della lunghezza della costa della Bretagna. A seconda della scala di misura (la lunghezza del "righello" utilizzato), la lunghezza totale misurata cambia. Questo perché a scale più piccole si possono cogliere dettagli e anfratti che vengono ignorati a scale più grandi. Se si osserva una foto di una roccia senza alcun riferimento dimensionale, è impossibile determinarne la grandezza reale. Zoomando sull'immagine, si continuano a osservare strutture statisticamente simili a quelle visibili a una scala più grande. L'invarianza di scala è il segnale matematico di questa proprietà.